%% kapa: exam=1 round=30 year=2019 period=1 form=A subject=civil_law subject_ko=민법 seq=24 number=24
%% source: Q-Net 감정평가사 기출문제 (cst00309) artlSeq=5205205 file=extracted/1008/5205205/2019년 제30회 감정평가사 1차시험 문제지/2019년 제30회 감정평가사 1차시험 문제지/2019년도 제30회 감정평가사 1차 1교시 A형.hwp
%% conversion: kordoc 4.24.51 -> tools/gen1.py (auto); review=proofread-sonnet; flags=none
\begin{question}{24}
무효등기의 유용에 관한 설명으로 옳지 않은 것은? (다툼이 있으면 판례에 따름)
\begin{choices}
  \item 무효등기의 유용에 관한 합의 내지 추인은 묵시적으로도 이루어질 수 있다.
  \item 실질관계의 소멸로 무효로 된 등기의 유용은 그 등기를 유용하기로 하는 합의가 이루어지기 전에 등기상 이해관계가 있는 제3자가 생기지 않은 경우에는 허용된다.
  \item 유용할 수 있는 등기에는 가등기도 포함된다.
  \item 기존건물이 전부 멸실된 후 그곳에 새로이 건축한 건물의 물권변동에 관한 등기를 위해 멸실된 건물의 등기를 유용할 수 있다.
  \item 무효인 등기를 유용하기로 한 약정을 하더라도, 무효의 등기가 있은 때로 소급하여 유효한 등기로 전환될 수 없다.
\end{choices}
\end{question}
\answer{4}
